\documentclass[aspectratio=169]{beamer}
\usetheme{metropolis}
\usepackage{graphicx}
\usepackage{tikz}
\usepackage{amsmath}
\usepackage{listings}
\usepackage{xcolor}
\usepackage{booktabs}
\usepackage{emoji}

% Code listing setup
\lstset{
  basicstyle=\ttfamily\small,
  keywordstyle=\color{blue},
  commentstyle=\color{gray},
  stringstyle=\color{red},
  showstringspaces=false,
  breaklines=true,
  frame=single,
  language=Python
}

\title{Contextual Embeddings: ELMo, USE, BERT}
\subtitle{Lecture 9: Beyond Static Word Representations}
\author{LLM Course - Week 4}
\date{Winter 2026}

\begin{document}

% ============================================================
% TITLE SLIDE
% ============================================================
\begin{frame}
\titlepage
\end{frame}

% ============================================================
% OVERVIEW
% ============================================================
\begin{frame}{Today's Lecture 📋}
\large
\begin{enumerate}
    \item 🔄 \textbf{From Static to Contextual}
    \item 🧬 \textbf{Language Models as Feature Extractors}
    \item 🔮 \textbf{ELMo: Embeddings from Language Models}
    \item 🌐 \textbf{Universal Sentence Encoder}
    \item 🎭 \textbf{BERT: Bidirectional Transformers}
    \item 📊 \textbf{Comparison \& Applications}
\end{enumerate}

\vspace{1em}
\textit{Goal: Understand how context transforms word representation}
\end{frame}

% ============================================================
% THE POLYSEMY PROBLEM
% ============================================================
\begin{frame}{The Polysemy Problem Revisited 🤔}
\textbf{Recall: Static embeddings assign ONE vector per word}

\vspace{1em}

\begin{exampleblock}{Example: "bank"}
\begin{enumerate}
    \item "I deposited money at the \textbf{bank}" \hfill (financial institution)
    \item "We sat by the river \textbf{bank}" \hfill (riverside)
    \item "The plane will \textbf{bank} left" \hfill (tilt/turn)
\end{enumerate}

\vspace{0.5em}
\textbf{Word2Vec/GloVe:} All three get \alert{the same vector} ✗

\textbf{Contextual embeddings:} Each gets a \alert{different vector} based on context ✓
\end{exampleblock}

\vspace{1em}

\begin{center}
\begin{tikzpicture}[scale=0.9]
    % Static
    \node at (0, 2) {\textbf{Static:}};
    \draw[fill=blue!30] (0,0.5) circle (0.3);
    \node at (0, 0) {bank};

    % Contextual
    \node at (5, 2) {\textbf{Contextual:}};
    \draw[fill=red!30] (3.5,0.5) circle (0.3);
    \node at (3.5, 0) {bank$_1$};

    \draw[fill=green!30] (5,0.5) circle (0.3);
    \node at (5, 0) {bank$_2$};

    \draw[fill=orange!30] (6.5,0.5) circle (0.3);
    \node at (6.5, 0) {bank$_3$};

    % Arrow
    \draw[->, very thick] (1, 1) -- (2.5, 1);
    \node at (1.75, 1.5) {Context};
\end{tikzpicture}
\end{center}
\end{frame}

% ============================================================
% STATIC VS CONTEXTUAL
% ============================================================
\begin{frame}[fragile]{Static vs. Contextual Embeddings 🔄}
\begin{columns}[T]
\begin{column}{0.48\textwidth}
\textbf{Static (Word2Vec, GloVe, FastText):}

\vspace{0.5em}
\begin{lstlisting}[language=Python]
# Same vector every time
model = Word2Vec(...)
vec1 = model['bank']
vec2 = model['bank']

assert vec1 == vec2  # True!
\end{lstlisting}

\vspace{0.5em}
\textbf{Characteristics:}
\begin{itemize}
    \item One vector per word type
    \item Context-independent
    \item Fast lookup (dictionary)
    \item Fixed after training
    \item Polysemy conflation
\end{itemize}
\end{column}

\begin{column}{0.48\textwidth}
\textbf{Contextual (ELMo, BERT):}

\vspace{0.5em}
\begin{lstlisting}[language=Python]
# Different vector per occurrence
model = BertModel.from_pretrained('bert-base')

sent1 = "river bank"
sent2 = "money bank"

vec1 = get_embedding(model, sent1, 'bank')
vec2 = get_embedding(model, sent2, 'bank')

assert vec1 != vec2  # True!
\end{lstlisting}

\vspace{0.5em}
\textbf{Characteristics:}
\begin{itemize}
    \item Different vector per occurrence
    \item Context-dependent
    \item Requires forward pass
    \item Dynamic representations
    \item Handles polysemy naturally
\end{itemize}
\end{column}
\end{columns}
\end{frame}

% ============================================================
% LANGUAGE MODELS
% ============================================================
\begin{frame}{Language Models as Feature Extractors 🧬}
\textbf{Key Insight:} Train a language model, use its internal states as embeddings

\vspace{1em}

\textbf{Language Modeling Task:}
\begin{center}
Predict the next word given previous words

\vspace{0.5em}
$P(w_t | w_1, w_2, ..., w_{t-1})$
\end{center}

\vspace{1em}

\begin{columns}[T]
\begin{column}{0.48\textwidth}
\textbf{Why This Works:}
\begin{itemize}
    \item To predict next word, model must understand:
    \begin{itemize}
        \item Syntax (grammar)
        \item Semantics (meaning)
        \item Context (discourse)
    \end{itemize}
    \item Hidden states encode this understanding
    \item Can extract and use as features!
\end{itemize}
\end{column}

\begin{column}{0.48\textwidth}
\begin{center}
\begin{tikzpicture}[scale=0.7]
    % Input
    \foreach \x/\word in {0/The, 1.5/cat, 3/sat, 4.5/on}
        \node at (\x, 4) {\word};

    % LSTM layers
    \foreach \x in {0, 1.5, 3, 4.5}
        \draw[fill=blue!20] (\x-0.3, 2.5) rectangle (\x+0.3, 3.2);

    \foreach \x in {0, 1.5, 3, 4.5}
        \draw[fill=green!20] (\x-0.3, 1.5) rectangle (\x+0.3, 2.2);

    % Arrows
    \foreach \x in {0, 1.5, 3, 4.5} {
        \draw[->] (\x, 3.7) -- (\x, 3.2);
        \draw[->] (\x, 2.5) -- (\x, 2.2);
    }

    % Predictions
    \foreach \x/\word in {0/cat, 1.5/sat, 3/on, 4.5/the}
        \node[font=\small] at (\x, 0.8) {\word};

    \foreach \x in {0, 1.5, 3, 4.5}
        \draw[->] (\x, 1.5) -- (\x, 1.1);

    % Label
    \node at (6, 2.85) {\alert{Extract these!}};
    \draw[->, red, thick] (5.5, 2.85) -- (5, 2.85);
\end{tikzpicture}
\end{center}

\vspace{0.5em}
\textit{Hidden states contain rich contextual information}
\end{column}
\end{columns}
\end{frame}

% ============================================================
% ELMO INTRO
% ============================================================
\begin{frame}{ELMo: Embeddings from Language Models 🔮}
\textbf{The first widely-adopted contextual embedding (2018)}

\vspace{1em}

\begin{columns}[T]
\begin{column}{0.48\textwidth}
\textbf{Key Ideas:}
\begin{enumerate}
    \item Train deep \alert{bidirectional LSTM} language model
    \item Use \alert{all layer} activations
    \item Weighted combination per task
    \item Pre-train on large corpus
\end{enumerate}

\vspace{0.5em}
\textbf{Architecture:}
\begin{itemize}
    \item 2-layer biLSTM
    \item Forward LM: $P(w_t | w_1...w_{t-1})$
    \item Backward LM: $P(w_t | w_{t+1}...w_n)$
    \item Character-based input (handles OOV!)
\end{itemize}
\end{column}

\begin{column}{0.48\textwidth}
\begin{center}
\begin{tikzpicture}[scale=0.6]
    % Input
    \node at (0, 0) {bank};

    % Character CNN
    \draw[fill=purple!20] (-0.8, 0.5) rectangle (0.8, 1.2);
    \node[font=\tiny] at (0, 0.85) {Char CNN};

    % Forward LSTM
    \draw[fill=blue!20] (-1.5, 2) rectangle (-0.3, 2.7);
    \node[font=\tiny] at (-0.9, 2.35) {LSTM$\rightarrow$};
    \draw[fill=blue!20] (-1.5, 3.2) rectangle (-0.3, 3.9);
    \node[font=\tiny] at (-0.9, 3.55) {LSTM$\rightarrow$};

    % Backward LSTM
    \draw[fill=red!20] (0.3, 2) rectangle (1.5, 2.7);
    \node[font=\tiny] at (0.9, 2.35) {LSTM$\leftarrow$};
    \draw[fill=red!20] (0.3, 3.2) rectangle (1.5, 3.9);
    \node[font=\tiny] at (0.9, 3.55) {LSTM$\leftarrow$};

    % Arrows up
    \draw[->] (0, 1.2) -- (-0.9, 2);
    \draw[->] (0, 1.2) -- (0.9, 2);
    \draw[->] (-0.9, 2.7) -- (-0.9, 3.2);
    \draw[->] (0.9, 2.7) -- (0.9, 3.2);

    % ELMo output
    \draw[fill=green!20] (-1, 4.5) rectangle (1, 5.2);
    \node[font=\small] at (0, 4.85) {ELMo};

    \draw[->] (-0.9, 3.9) -- (0, 4.5);
    \draw[->] (0.9, 3.9) -- (0, 4.5);

    % Labels
    \node[font=\tiny] at (-2.5, 2.35) {Layer 1};
    \node[font=\tiny] at (-2.5, 3.55) {Layer 2};
\end{tikzpicture}
\end{center}

\vspace{0.5em}
\textbf{ELMo representation:}
\begin{align*}
\text{ELMo}_k^{task} = \gamma^{task} \sum_{j=0}^{L} s_j^{task} \mathbf{h}_{k,j}
\end{align*}

Learned task-specific weights!
\end{column}
\end{columns}

\vspace{0.5em}
\small
\textit{Reference: Peters et al. (2018). "Deep contextualized word representations"}
\end{frame}

% ============================================================
% ELMO DETAILS
% ============================================================
\begin{frame}{ELMo: How It Works 🔍}
\textbf{Training:}

\vspace{0.5em}
\begin{enumerate}
    \item Pre-train on large corpus (1B Word Benchmark)
    \item Minimize combined forward \& backward LM loss:
    \begin{align*}
    \mathcal{L} = \sum_{k=1}^N \left(\log P(w_k | w_1...w_{k-1}) + \log P(w_k | w_{k+1}...w_N)\right)
    \end{align*}
    \item Each position gets representation from all layers
\end{enumerate}

\vspace{1em}

\textbf{Usage (downstream tasks):}

\vspace{0.5em}
\begin{enumerate}
    \item Freeze ELMo weights
    \item For each token, compute:
    \begin{itemize}
        \item Character-based representation
        \item 2 forward LSTM hidden states
        \item 2 backward LSTM hidden states
        \item Total: 5 representations (char + 2×2 layers)
    \end{itemize}
    \item Learn task-specific weighted combination
    \item Concatenate with task model
\end{enumerate}

\vspace{0.5em}
\textbf{Impact:} Improved state-of-the-art on 6 NLP tasks!
\end{frame}

% ============================================================
% ELMO CODE
% ============================================================
\begin{frame}[fragile]{ELMo in Practice 💻}
\begin{lstlisting}[language=Python]
from allennlp.modules.elmo import Elmo, batch_to_ids

# Initialize ELMo
options_file = "https://s3-us-west-2.amazonaws.com/allennlp/models/elmo/2x4096_512_2048cnn_2xhighway/elmo_2x4096_512_2048cnn_2xhighway_options.json"
weight_file = "https://s3-us-west-2.amazonaws.com/allennlp/models/elmo/2x4096_512_2048cnn_2xhighway/elmo_2x4096_512_2048cnn_2xhighway_weights.hdf5"

elmo = Elmo(options_file, weight_file, 2, dropout=0)

# Prepare sentences
sentences = [
    ['I', 'deposited', 'money', 'at', 'the', 'bank'],
    ['We', 'sat', 'by', 'the', 'river', 'bank']
]

# Convert to character ids
character_ids = batch_to_ids(sentences)

# Get embeddings
embeddings = elmo(character_ids)

# embeddings['elmo_representations'] contains:
# - List of 2 tensors (one per layer)
# - Shape: [batch_size, seq_len, 1024]

# Different vectors for "bank"!
bank1 = embeddings['elmo_representations'][0][0, 5, :]  # first sentence
bank2 = embeddings['elmo_representations'][0][1, 5, :]  # second sentence

# Cosine similarity will be lower than for static embeddings
\end{lstlisting}
\end{frame}

% ============================================================
% UNIVERSAL SENTENCE ENCODER
% ============================================================
\begin{frame}{Universal Sentence Encoder (USE) 🌐}
\textbf{Sentence-level embeddings for semantic similarity}

\vspace{1em}

\begin{columns}[T]
\begin{column}{0.48\textwidth}
\textbf{Motivation:}
\begin{itemize}
    \item Word embeddings: good for words
    \item But what about sentences?
    \item Average of word vectors? Too simple!
    \item Need compositionality
\end{itemize}

\vspace{0.5em}
\textbf{Two Variants:}

\vspace{0.3em}
\textbf{1. Transformer-based:}
\begin{itemize}
    \item Higher accuracy
    \item Slower (compute intensive)
    \item Better for quality
\end{itemize}

\vspace{0.3em}
\textbf{2. Deep Averaging Network (DAN):}
\begin{itemize}
    \item Lower accuracy
    \item Much faster
    \item Better for scale
\end{itemize}
\end{column}

\begin{column}{0.48\textwidth}
\textbf{Training Objectives:}
\begin{enumerate}
    \item Unsupervised: Skip-thought
    \item Supervised: SNLI (entailment)
    \item Multi-task learning
\end{enumerate}

\vspace{0.5em}
\textbf{Output:}
\begin{itemize}
    \item 512-dimensional vector
    \item Fixed-length (any sentence length)
    \item Optimized for similarity tasks
\end{itemize}

\vspace{0.5em}
\textbf{Use Cases:}
\begin{itemize}
    \item Semantic search
    \item Question answering
    \item Text classification
    \item Clustering
    \item Duplicate detection
\end{itemize}
\end{column}
\end{columns}

\vspace{0.5em}
\small
\textit{Reference: Cer et al. (2018). "Universal Sentence Encoder"}
\end{frame}

% ============================================================
% USE CODE
% ============================================================
\begin{frame}[fragile]{Universal Sentence Encoder in Practice 💻}
\begin{lstlisting}[language=Python]
import tensorflow_hub as hub
import numpy as np

# Load model
embed = hub.load("https://tfhub.dev/google/universal-sentence-encoder/4")

# Example sentences
sentences = [
    "The cat sat on the mat.",
    "A feline rested on the rug.",
    "The dog ran in the park.",
    "I love machine learning."
]

# Generate embeddings
embeddings = embed(sentences)

# Shape: [4, 512]
print(embeddings.shape)

# Compute similarity
from sklearn.metrics.pairwise import cosine_similarity

sim_matrix = cosine_similarity(embeddings)
print(sim_matrix)

# Sentences 1 and 2 should be very similar (paraphrases)
# Sentence 3 somewhat similar (animals)
# Sentence 4 dissimilar

# Use for semantic search
query = "cat on mat"
query_embedding = embed([query])
similarities = cosine_similarity(query_embedding, embeddings)[0]
most_similar_idx = np.argmax(similarities)
print(f"Most similar: {sentences[most_similar_idx]}")
\end{lstlisting}
\end{frame}

% ============================================================
% BERT INTRO
% ============================================================
\begin{frame}{BERT: Bidirectional Encoder Representations 🎭}
\begin{center}
\LARGE
\textbf{The model that changed everything (2018)}
\end{center}

\vspace{1em}

\begin{columns}[T]
\begin{column}{0.48\textwidth}
\textbf{Key Innovations:}
\begin{enumerate}
    \item \alert{Bidirectional} context (not just left-to-right)
    \item \alert{Transformer} architecture (attention)
    \item \alert{Masked language model} pre-training
    \item \alert{Next sentence prediction}
    \item Deeply bidirectional
\end{enumerate}

\vspace{0.5em}
\textbf{Impact:}
\begin{itemize}
    \item SOTA on 11 NLP tasks
    \item Sparked the "BERT-era"
    \item 1000+ variants (RoBERTa, ALBERT, DistilBERT, ...)
    \item Foundation for modern LLMs
\end{itemize}
\end{column}

\begin{column}{0.48\textwidth}
\textbf{Architecture Sizes:}

\vspace{0.5em}
\begin{tabular}{lcc}
\toprule
\textbf{Model} & \textbf{Layers} & \textbf{Hidden} \\
\midrule
BERT-Base & 12 & 768 \\
BERT-Large & 24 & 1024 \\
\bottomrule
\end{tabular}

\vspace{1em}
\textbf{Training Data:}
\begin{itemize}
    \item BooksCorpus (800M words)
    \item English Wikipedia (2.5B words)
    \item Total: 3.3B words
\end{itemize}

\vspace{0.5em}
\textbf{Training Time:}
\begin{itemize}
    \item 4 days on 64 TPU chips
    \item Or 4 weeks on 8 GPUs
\end{itemize}
\end{column}
\end{columns}

\vspace{0.5em}
\small
\textit{Reference: Devlin et al. (2018). "BERT: Pre-training of Deep Bidirectional Transformers"}
\end{frame}

% ============================================================
% MASKED LM
% ============================================================
\begin{frame}{Masked Language Modeling (MLM) 🎭}
\textbf{BERT's key training innovation}

\vspace{1em}

\textbf{The Problem with Traditional LM:}
\begin{itemize}
    \item Left-to-right: Only sees previous words
    \item Right-to-left: Only sees next words
    \item Want: See both directions simultaneously
    \item But: Can't just show the answer during training!
\end{itemize}

\vspace{1em}

\textbf{Solution: Mask some words, predict them}

\vspace{0.5em}

\begin{exampleblock}{Example}
\textbf{Original:} "The cat sat on the mat"

\textbf{Masked:} "The [MASK] sat on the mat"

\textbf{Task:} Predict [MASK] = "cat"
\end{exampleblock}

\vspace{1em}

\textbf{Training Procedure:}
\begin{enumerate}
    \item Randomly select 15\% of tokens
    \item Replace 80\% with [MASK], 10\% with random word, 10\% unchanged
    \item Predict original tokens
    \item Bidirectional context!
\end{enumerate}
\end{frame}

% ============================================================
% NSP
% ============================================================
\begin{frame}{Next Sentence Prediction (NSP) 🔗}
\textbf{Second pre-training task: Understand sentence relationships}

\vspace{1em}

\textbf{Task:} Given two sentences A and B, predict if B follows A in the text

\vspace{0.5em}

\begin{exampleblock}{Positive Example (IsNext)}
\textbf{Sentence A:} "The cat sat on the mat."

\textbf{Sentence B:} "It was sleeping peacefully."

\textbf{Label:} IsNext ✓
\end{exampleblock}

\vspace{0.5em}

\begin{exampleblock}{Negative Example (NotNext)}
\textbf{Sentence A:} "The cat sat on the mat."

\textbf{Sentence B:} "Machine learning is fascinating."

\textbf{Label:} NotNext ✗
\end{exampleblock}

\vspace{1em}

\textbf{Why This Helps:}
\begin{itemize}
    \item Captures discourse relationships
    \item Useful for QA, NLI, etc.
    \item Models sentence-level coherence
    \item Note: Later research (RoBERTa) found NSP less important than MLM
\end{itemize}
\end{frame}

% ============================================================
% BERT ARCHITECTURE
% ============================================================
\begin{frame}{BERT Architecture 🏗️}
\begin{center}
\begin{tikzpicture}[scale=0.7]
    % Input
    \node at (-3, 0) {Input:};
    \foreach \x/\word in {0/[CLS], 1.5/The, 3/cat, 4.5/[MASK], 6/on, 7.5/mat, 9/[SEP]}
        \node[font=\small] at (\x, 0) {\word};

    % Token Embeddings
    \foreach \x in {0, 1.5, 3, 4.5, 6, 7.5, 9}
        \draw[fill=blue!20] (\x-0.4, 0.8) rectangle (\x+0.4, 1.3);
    \node at (-3, 1.05) {Token:};

    % Segment Embeddings
    \foreach \x in {0, 1.5, 3, 4.5, 6, 7.5}
        \draw[fill=green!20] (\x-0.4, 1.5) rectangle (\x+0.4, 2);
    \draw[fill=orange!20] (9-0.4, 1.5) rectangle (9+0.4, 2);
    \node at (-3, 1.75) {Segment:};

    % Position Embeddings
    \foreach \x/\pos in {0/0, 1.5/1, 3/2, 4.5/3, 6/4, 7.5/5, 9/6}
        \draw[fill=red!20] (\x-0.4, 2.2) rectangle (\x+0.4, 2.7);
    \node at (-3, 2.45) {Position:};

    % Sum
    \node at (3, 3.2) {$+$};

    % Transformer layers
    \draw[fill=purple!20, rounded corners] (-0.5, 3.7) rectangle (9.5, 4.5);
    \node at (4.5, 4.1) {Transformer Layer 1};

    \draw[fill=purple!20, rounded corners] (-0.5, 4.7) rectangle (9.5, 5.5);
    \node at (4.5, 5.1) {Transformer Layer 2};

    \node at (4.5, 6) {$\vdots$};

    \draw[fill=purple!20, rounded corners] (-0.5, 6.3) rectangle (9.5, 7.1);
    \node at (4.5, 6.7) {Transformer Layer 12};

    % Output
    \foreach \x in {0, 1.5, 3, 4.5, 6, 7.5, 9}
        \draw[fill=yellow!30] (\x-0.4, 7.5) rectangle (\x+0.4, 8);
    \node at (-3, 7.75) {Output:};

    % Task heads
    \draw[->] (0, 8) -- (0, 8.5);
    \node[font=\small] at (0, 9) {Classification};

    \draw[->] (4.5, 8) -- (4.5, 8.5);
    \node[font=\small] at (4.5, 9) {Predict: sat};
\end{tikzpicture}
\end{center}
\end{frame}

% ============================================================
% BERT CODE
% ============================================================
\begin{frame}[fragile]{BERT in Practice 💻}
\begin{lstlisting}[language=Python]
from transformers import BertTokenizer, BertModel
import torch

# Load pre-trained BERT
tokenizer = BertTokenizer.from_pretrained('bert-base-uncased')
model = BertModel.from_pretrained('bert-base-uncased')

# Example sentences with "bank"
sent1 = "I deposited money at the bank"
sent2 = "We sat by the river bank"

# Tokenize
tokens1 = tokenizer(sent1, return_tensors='pt')
tokens2 = tokenizer(sent2, return_tensors='pt')

# Get embeddings
with torch.no_grad():
    output1 = model(**tokens1)
    output2 = model(**tokens2)

# Last hidden state: [batch_size, seq_len, hidden_size]
embeddings1 = output1.last_hidden_state
embeddings2 = output2.last_hidden_state

# Extract "bank" embedding (position varies)
# tokens1: [CLS] i deposited money at the bank [SEP]
bank1_embedding = embeddings1[0, 6, :]  # 768-dim vector

# tokens2: [CLS] we sat by the river bank [SEP]
bank2_embedding = embeddings2[0, 6, :]  # 768-dim vector

# Different vectors for "bank"!
from torch.nn.functional import cosine_similarity
sim = cosine_similarity(bank1_embedding, bank2_embedding, dim=0)
print(f"Similarity: {sim:.3f}")  # Lower than with static embeddings
\end{lstlisting}
\end{frame}

% ============================================================
% FINE-TUNING
% ============================================================
\begin{frame}[fragile]{Fine-tuning BERT 🎯}
\textbf{Two ways to use BERT:}

\vspace{1em}

\begin{columns}[T]
\begin{column}{0.48\textwidth}
\textbf{1. Feature Extraction:}
\begin{itemize}
    \item Freeze BERT weights
    \item Use embeddings as features
    \item Train classifier on top
    \item Faster, less data needed
\end{itemize}

\vspace{0.5em}
\begin{lstlisting}[language=Python]
# Freeze BERT
for param in bert_model.parameters():
    param.requires_grad = False

# Add classifier
classifier = nn.Linear(768, num_classes)

# Train only classifier
optimizer = Adam(classifier.parameters())
\end{lstlisting}
\end{column}

\begin{column}{0.48\textwidth}
\textbf{2. Fine-tuning:}
\begin{itemize}
    \item Update BERT weights
    \item Add task-specific head
    \item Train end-to-end
    \item Better performance, more data needed
\end{itemize}

\vspace{0.5em}
\begin{lstlisting}[language=Python]
# Keep BERT trainable
bert_model = BertModel.from_pretrained(
    'bert-base-uncased'
)

# Add classifier
classifier = nn.Linear(768, num_classes)

# Train everything
optimizer = Adam(
    list(bert_model.parameters()) +
    list(classifier.parameters()),
    lr=2e-5  # Small learning rate!
)
\end{lstlisting}
\end{column}
\end{columns}

\vspace{1em}
\textbf{Best Practice:} Fine-tune with small learning rate (1e-5 to 5e-5) for few epochs (2-4)
\end{frame}

% ============================================================
% COMPARISON
% ============================================================
\begin{frame}{Contextual Embeddings Comparison 📊}
\begin{center}
\begin{tabular}{lcccc}
\toprule
\textbf{Property} & \textbf{ELMo} & \textbf{USE} & \textbf{BERT} \\
\midrule
Architecture & BiLSTM & Transformer/DAN & Transformer \\
Pre-training & LM (forward+backward) & Multi-task & MLM + NSP \\
Bidirectional & Shallow & Yes & Deep \\
Granularity & Token & Sentence & Token \\
Hidden size & 1024 & 512 & 768/1024 \\
Parameters & 93M & 256M & 110M/340M \\
Speed & Medium & Fast & Slow \\
OOV handling & Characters & Subwords & WordPiece \\
Year & 2018 & 2018 & 2018 \\
\bottomrule
\end{tabular}
\end{center}

\vspace{1em}

\begin{block}{Recommendations}
\begin{itemize}
    \item \textbf{ELMo:} Legacy systems, character-aware needs
    \item \textbf{USE:} Sentence similarity, semantic search, fast inference
    \item \textbf{BERT:} State-of-the-art quality, most NLP tasks (now superseded by larger models)
\end{itemize}
\end{block}
\end{frame}

% ============================================================
% IMPACT
% ============================================================
\begin{frame}{Impact on NLP 🌟}
\textbf{BERT revolutionized NLP:}

\vspace{1em}

\begin{columns}[T]
\begin{column}{0.48\textwidth}
\textbf{Before BERT (pre-2018):}
\begin{itemize}
    \item Task-specific architectures
    \item Train from scratch
    \item Static embeddings (Word2Vec, GloVe)
    \item Limited transfer learning
    \item Moderate performance
\end{itemize}

\vspace{1em}
\textbf{Tasks that improved:}
\begin{itemize}
    \item Question Answering (+1.5 F1 on SQuAD)
    \item NER (+0.3 F1)
    \item Sentiment Analysis (+2\% accuracy)
    \item NLI (+4\% accuracy)
    \item Many others!
\end{itemize}
\end{column}

\begin{column}{0.48\textwidth}
\textbf{After BERT (post-2018):}
\begin{itemize}
    \item Pre-train, then fine-tune
    \item Transfer learning standard
    \item Contextual embeddings
    \item Massive pre-trained models
    \item State-of-the-art results
\end{itemize}

\vspace{1em}
\textbf{BERT Variants:}
\begin{itemize}
    \item RoBERTa: Better training
    \item ALBERT: Parameter sharing
    \item DistilBERT: Smaller, faster
    \item ELECTRA: Different pre-training
    \item DeBERTa: Disentangled attention
    \item And 100+ more!
\end{itemize}
\end{column}
\end{columns}

\vspace{1em}
\begin{center}
\alert{BERT established the "pre-train + fine-tune" paradigm that dominates NLP today}
\end{center}
\end{frame}

% ============================================================
% APPLICATIONS
% ============================================================
\begin{frame}{Real-World Applications 🚀}
\begin{enumerate}
    \item \textbf{Search \& Information Retrieval:}
    \begin{itemize}
        \item Google Search uses BERT for query understanding
        \item Semantic search engines
        \item Document ranking
    \end{itemize}

    \item \textbf{Question Answering:}
    \begin{itemize}
        \item Reading comprehension (SQuAD)
        \item Customer support chatbots
        \item FAQ systems
    \end{itemize}

    \item \textbf{Text Classification:}
    \begin{itemize}
        \item Sentiment analysis
        \item Intent detection
        \item Content moderation
    \end{itemize}

    \item \textbf{Named Entity Recognition:}
    \begin{itemize}
        \item Extract people, places, organizations
        \item Medical entity extraction
        \item Legal document processing
    \end{itemize}

    \item \textbf{Semantic Similarity:}
    \begin{itemize}
        \item Duplicate detection
        \item Paraphrase identification
        \item Recommendation systems
    \end{itemize}
\end{enumerate}
\end{frame}

% ============================================================
% DISCUSSION
% ============================================================
\begin{frame}{Discussion Question 💬}
\begin{center}
\Large
\textbf{Do contextual embeddings truly "understand" language?}
\end{center}

\vspace{1em}

\textbf{Consider:}
\begin{itemize}
    \item BERT can distinguish "bank" (financial) from "bank" (riverside)
    \item It achieves human-level performance on many benchmarks
    \item But it's trained only on text co-occurrence patterns
\end{itemize}

\vspace{1em}

\begin{columns}[T]
\begin{column}{0.48\textwidth}
\textbf{Arguments For:}
\begin{itemize}
    \item Captures complex semantic relationships
    \item Generalizes to new contexts
    \item Emergent linguistic capabilities
    \item Handles compositional meaning
\end{itemize}
\end{column}

\begin{column}{0.48\textwidth}
\textbf{Arguments Against:}
\begin{itemize}
    \item No grounding in physical world
    \item No common sense reasoning
    \item Exploits statistical shortcuts
    \item Brittle to adversarial examples
    \item "Stochastic parrots"?
\end{itemize}
\end{column}
\end{columns}

\vspace{1em}

\begin{alertblock}{The Grounding Problem Persists}
Even with contextual embeddings, we still lack true grounding in experience, perception, and embodied cognition.
\end{alertblock}

\small
\textit{Reference: Bender \& Koller (2020). "Climbing towards NLU: On Meaning, Form, and Understanding"}
\end{frame}

% ============================================================
% PRACTICAL TIPS
% ============================================================
\begin{frame}{Practical Tips 💡}
\begin{enumerate}
    \item \textbf{Choosing a Model:}
    \begin{itemize}
        \item BERT-base: Good balance, 110M params
        \item DistilBERT: 40\% smaller, 60\% faster, 97\% performance
        \item RoBERTa: Better than BERT, longer training
        \item Domain-specific: BioBERT, SciBERT, FinBERT, etc.
    \end{itemize}

    \item \textbf{Fine-tuning Best Practices:}
    \begin{itemize}
        \item Small learning rate (2e-5 typical)
        \item Few epochs (2-4)
        \item Batch size: 16 or 32
        \item Warm-up steps
        \item Gradient clipping
    \end{itemize}

    \item \textbf{Computational Considerations:}
    \begin{itemize}
        \item BERT-base: ~110M params, 512 max tokens
        \item Needs GPU (1-4 GB VRAM minimum)
        \item Batching for efficiency
        \item Consider DistilBERT for production
    \end{itemize}

    \item \textbf{Using HuggingFace:}
    \begin{itemize}
        \item Easy access to 1000+ pre-trained models
        \item Standardized API
        \item Good documentation and community
    \end{itemize}
\end{enumerate}
\end{frame}

% ============================================================
% SUMMARY
% ============================================================
\begin{frame}{Summary 🎯}
\textbf{What we learned today:}

\begin{enumerate}
    \item \textbf{Contextual vs. Static:} Different vectors per occurrence

    \item \textbf{ELMo (2018):}
    \begin{itemize}
        \item BiLSTM language models
        \item Character-based, handles OOV
        \item Task-specific weighting
    \end{itemize}

    \item \textbf{Universal Sentence Encoder (2018):}
    \begin{itemize}
        \item Sentence-level embeddings
        \item Two variants: Transformer \& DAN
        \item Optimized for semantic similarity
    \end{itemize}

    \item \textbf{BERT (2018):}
    \begin{itemize}
        \item Masked language modeling
        \item Deep bidirectional transformers
        \item Pre-train + fine-tune paradigm
        \item Revolutionized NLP
    \end{itemize}

    \item \textbf{Impact:} Established modern transfer learning in NLP

    \item \textbf{Next:} Dimensionality reduction techniques for visualization
\end{enumerate}
\end{frame}

% ============================================================
% REFERENCES
% ============================================================
\begin{frame}{Key References 📚}
\small
\textbf{Foundational Papers:}
\begin{itemize}
    \item Peters et al. (2018). "Deep contextualized word representations" (ELMo)
    \item Cer et al. (2018). "Universal Sentence Encoder"
    \item Devlin et al. (2018). "BERT: Pre-training of Deep Bidirectional Transformers for Language Understanding"
\end{itemize}

\vspace{0.5em}
\textbf{BERT Variants:}
\begin{itemize}
    \item Liu et al. (2019). "RoBERTa: A Robustly Optimized BERT Pretraining Approach"
    \item Sanh et al. (2019). "DistilBERT, a distilled version of BERT"
    \item Lan et al. (2019). "ALBERT: A Lite BERT for Self-supervised Learning"
\end{itemize}

\vspace{0.5em}
\textbf{Critical Perspectives:}
\begin{itemize}
    \item Bender \& Koller (2020). "Climbing towards NLU: On Meaning, Form, and Understanding"
    \item Bender et al. (2021). "On the Dangers of Stochastic Parrots"
\end{itemize}

\vspace{0.5em}
\textbf{Resources:}
\begin{itemize}
    \item HuggingFace Transformers: \url{https://huggingface.co/transformers/}
    \item BERT Paper: \url{https://arxiv.org/abs/1810.04805}
\end{itemize}
\end{frame}

% ============================================================
% QUESTIONS
% ============================================================
\begin{frame}{Questions? 🙋}
\begin{center}
\LARGE
\textbf{Next Lecture:}

\vspace{1em}
Dimensionality Reduction: PCA, t-SNE, UMAP

\vspace{2em}
\textit{Visualizing high-dimensional embeddings!}
\end{center}
\end{frame}

\end{document}
